% Generated by analysis/v79_cond_audit.py; frozen predictions unchanged.
\begin{table}[H]
\setbox2=\vbox{\begin{minipage}{\linewidth}
\centering\small
\setlength{\tabcolsep}{4pt}
\caption{The V79 audit examines capability conditioning. Variant B uses unrestricted signed scales. Pruning confirmation counts the two identical Pythia-2.8B revisions as one state. Positive $\Delta=B-A$ is the MAE reduction from A, in nats. Rows identify compression families above and capabilities below; columns give dimensionless scales and squared correlations above, and error comparisons below. This is a \mbox{post-hoc} audit. A denotes a separate response per capability; B denotes a shared response with a fitted capability scale. CI means confidence interval; the numbered records identify the original audits and confirmations. MAE denotes mean absolute error. CI denotes confidence interval.}
\label{tab:cond_audit}
\setbox0=\hbox{
\begin{tabular}{lrrr}
\toprule
\shortstack[l]{Response family or\\diagnostic} & Math & Code & QA \\
\midrule
Pruning & 1.287773 & 1.270559 & 0.327894 \\
Quantization & 1.069770 & 1.431306 & 0.556505 \\
Distillation & 0.716184 & 0.734226 & 1.549590 \\
\midrule
Pruning anchor $R^2$ & 0.9382 & 0.9739 & 0.4894 \\
\bottomrule
\end{tabular}
}\ifdim\wd0>\linewidth\resizebox{\linewidth}{!}{\box0}\else\box0\fi
\par\medskip
\resizebox{\linewidth}{!}{%
\setbox0=\hbox{
\begin{tabular}{lrrll}
\toprule
Capability & A, 4 states & B, 4 states & $\Delta$, 5 labels [95\% CI] & $\Delta$, 4 states [95\% CI] \\
\midrule
Math & 0.2279 & 0.2328 & \shortstack[l]{0.0190 [-0.0291,\\0.0690]} & \shortstack[l]{0.0048 [-0.0495,\\0.0591]} \\
Code & 0.1736 & 0.1776 & \shortstack[l]{0.0092 [-0.0116,\\0.0273]} & \shortstack[l]{0.0040 [-0.0181,\\0.0228]} \\
QA & 0.1818 & 0.3500 & \shortstack[l]{0.1899 [0.1252,\\0.2545]} & \shortstack[l]{0.1683 [0.1054,\\0.2434]} \\
Capability mean & 0.1944 & 0.2535 & \shortstack[l]{0.0727 [0.0300,\\0.1154]} & \shortstack[l]{0.0590 [0.0206,\\0.1025]} \\
\bottomrule
\end{tabular}
}\ifdim\wd0>\linewidth\resizebox{\linewidth}{!}{\box0}\else\box0\fi
}
\par\smallskip
\begin{minipage}{0.99\linewidth}\footnotesize All fitted scales, including stored sensitivities, are positive, although negatives are permitted. $R^2=1-\sum_d(m_c-s_ch)^2/\sum_d(m_c-\bar m_c)^2$ uses seven development median anchors and no fitted intercept. QA is positive at $d=0.55$ and negative at $d=0.60,0.65,0.70,0.75,0.80,0.90$; Math and Code medians remain positive. The shares of sources opposite the arithmetic-mean QA direction at those seven densities are $1/8,10/16,7/8,12/16,0/4,14/16,10/16$, respectively. The gain primarily reflects non-proportional aggregate QA shape; neither source-free variant learns source-specific direction. It is not caused by a nonnegative-scale constraint. The four-state comparison gives each duplicate cell one vote (12 cells, 36 capability rows); 5,000 paired state-bootstrap draws, seed 0, fixed fits. Intervals are descriptive with few states. QA and capability-mean intervals remain positive and Math and Code intervals include zero; these results reproduce V76's unique-outcome sensitivity. Keeping all 45 rows but merging bootstrap labels preserves the original point estimates (see accompanying summary). V72 freeze/compare also count the two aliases separately; no other such files matched both tags. \textbf{Distillation A$=$B is algebraic for the single-coefficient forms ($a_c=a s_c$, $a\ne0$) and supports no conclusion about conditioning.} \end{minipage}
\end{minipage}}
\centering\ifdim\dimexpr\ht2+\dp2\relax>.95\textheight\resizebox*{!}{.95\textheight}{\box2}\else\box2\fi
\end{table}
